\chapter{Penetration Testing}
\label{ch:chapter04}
"[P][...]enetration testing is a detailed procedure to unravel [...] open doors in a system or network." \cite{Lachkov_2022} The technique of penetration testing is a creative process of understanding IT-systems and their detailed structure and functionality. Finding and documenting vulnerabilities in IT-infrastructure helps to secure them and enhance the safety of these systems.

\ac{DAST} scanners test software for vulnerabilities by replaying known attack patterns. In comparison, penetration tests are more creative approaches. Testers adapt their testing plans based on system reactions and architectural details. Penetration tests simulate realistic attacks against a system under test, revealing in detail how real-world attackers could exploit existing vulnerabilities. This information is documented in a test report to inform IT infrastructure management and propose mitigation strategies. \cite{Lachkov_2022}

If a corporate vendor decides to conduct a penetration test on its system, this test can be performed either by internal employees or by a specialized third-party company. Either way the goal of finding exploitable vulnerabilities needs to be specified in detail. You can refer to the example of stealing a car. By a penetration test, a specialist can find out how hard it is to steal your car, but this appears not to be quite specific. Stealing a car usually refers to someone driving another's car without permission and not returning it. Does an unknown driving one's car without permission refers to stealing it when it is returned the next day? Does a teenager driving its parents car without parents permission refer to stealing it? The goal of finding out how hard it is to steal a car is simply not specific enough. For IT infrastructure a specific goal could for example be: Is an unauthorized party able to interfere with \ac{CIA} triad? \cite{Bishop_2007}

In the next step it is important to consider the knowledge the penetration tester has about the system under test. Potential attackers can allocate massive resources in the form of time and money if the exploitation seems worthwhile. Through providing penetration testers with information about the system under test they can compensate for the disparity in resources. \cite{Bishop_2007} Also penetration testers can be given partially internal access to simulate the scenario of already overcome first barriers by attackers. \cite{Bishop_2007} Penetration testing falls into three categories based on knowledge level: Black-Box, Grey-Box, and White-Box penetration testing. \cite{Rane_2024}

% more could be written about stages of penetration testing